\documentclass[10pt,a4paper]{amsart}
\usepackage{amsmath,amssymb,amsthm,graphicx}
\usepackage{verbatim}
\newenvironment{reference}{\comment}{\endcomment}
%\newenvironment{reference}{}{}
\newenvironment{slogan}{\comment}{\endcomment}
\newenvironment{history}{\comment}{\endcomment}
% For commutative diagrams we use Xy-pic
\usepackage[all]{xy}

% We use 2cell for 2-commutative diagrams.
\xyoption{2cell}
% Theorem environments.
%
\theoremstyle{plain}
\newtheorem{theorem}[subsection]{Theorem}
\newtheorem{proposition}[subsection]{Proposition}
\newtheorem{lemma}[subsection]{Lemma}

\theoremstyle{definition}
\newtheorem{definition}[subsection]{Definition}
\newtheorem{example}[subsection]{Example}
\newtheorem{exercise}[subsection]{Exercise}
\newtheorem{situation}[subsection]{Situation}

\theoremstyle{remark}
\newtheorem{remark}[subsection]{Remark}
\newtheorem{remarks}[subsection]{Remarks}

\numberwithin{equation}{subsection}

% Macros
%
\def\lim{\mathop{\mathrm{lim}}\nolimits}
\def\colim{\mathop{\mathrm{colim}}\nolimits}
\def\Spec{\mathop{\mathrm{Spec}}}
\def\Hom{\mathop{\mathrm{Hom}}\nolimits}
\def\Ext{\mathop{\mathrm{Ext}}\nolimits}
\def\SheafHom{\mathop{\mathcal{H}\!\mathit{om}}\nolimits}
\def\SheafExt{\mathop{\mathcal{E}\!\mathit{xt}}\nolimits}
\def\Sch{\mathit{Sch}}
\def\Mor{\mathop{\mathrm{Mor}}\nolimits}
\def\Ob{\mathop{\mathrm{Ob}}\nolimits}
\def\Sh{\mathop{\mathit{Sh}}\nolimits}
\def\NL{\mathop{N\!L}\nolimits}
\def\CH{\mathop{\mathrm{CH}}\nolimits}
\def\proetale{{pro\text{-}\acute{e}tale}}
\def\etale{{\acute{e}tale}}
\def\QCoh{\mathit{QCoh}}
\def\Ker{\mathop{\mathrm{Ker}}}
\def\Im{\mathop{\mathrm{Im}}}
\def\Coker{\mathop{\mathrm{Coker}}}
\def\Coim{\mathop{\mathrm{Coim}}}

% Boxtimes
%
\DeclareMathSymbol{\boxtimes}{\mathbin}{AMSa}{"02}

%
% Macros for moduli stacks/spaces
%
\def\QCohstack{\mathcal{QC}\!\mathit{oh}}
\def\Cohstack{\mathcal{C}\!\mathit{oh}}
\def\Spacesstack{\mathcal{S}\!\mathit{paces}}
\def\Quotfunctor{\mathrm{Quot}}
\def\Hilbfunctor{\mathrm{Hilb}}
\def\Curvesstack{\mathcal{C}\!\mathit{urves}}
\def\Polarizedstack{\mathcal{P}\!\mathit{olarized}}
\def\Complexesstack{\mathcal{C}\!\mathit{omplexes}}
% \Pic is the operator that assigns to X its picard group, usage \Pic(X)
% \Picardstack_{X/B} denotes the Picard stack of X over B
% \Picardfunctor_{X/B} denotes the Picard functor of X over B
\def\Pic{\mathop{\mathrm{Pic}}\nolimits}
\def\Picardstack{\mathcal{P}\!\mathit{ic}}
\def\Picardfunctor{\mathrm{Pic}}
\def\Deformationcategory{\mathcal{D}\!\mathit{ef}}
\usepackage{cleveref}
\crefname{originallemma}{original lemma}{original lemmas}
\Crefname{originallemma}{Original lemma}{Original lemmas}
\makeatletter
\newlabel{coherent-lemma-quasi-coherent-affine-cohomology-zero}{{01XB}{source}{}{sourcepointer-1}{}}
\newlabel{coherent-lemma-quasi-coherent-affine-cohomology-zero@cref}{{[originallemma][0][]01XB}{[1][0][]source}{}{}{}}
\makeatother
\crefname{originallemma}{original lemma}{original lemmas}
\Crefname{originallemma}{Original lemma}{Original lemmas}
\makeatletter
\newlabel{lemma-coherent-Noetherian}{{07UB}{source}{}{sourcepointer-2}{}}
\newlabel{lemma-coherent-Noetherian@cref}{{[originallemma][0][]07UB}{[1][0][]source}{}{}{}}
\makeatother
\crefname{originallemma}{original lemma}{original lemmas}
\Crefname{originallemma}{Original lemma}{Original lemmas}
\makeatletter
\newlabel{lemma-higher-direct-image}{{0720}{source}{}{sourcepointer-3}{}}
\newlabel{lemma-higher-direct-image@cref}{{[originallemma][0][]0720}{[1][0][]source}{}{}{}}
\makeatother
\crefname{originallemma}{original lemma}{original lemmas}
\Crefname{originallemma}{Original lemma}{Original lemmas}
\makeatletter
\newlabel{lemma-proper-pushforward-coherent}{{08AR}{source}{}{sourcepointer-4}{}}
\newlabel{lemma-proper-pushforward-coherent@cref}{{[originallemma][0][]08AR}{[1][0][]source}{}{}{}}
\makeatother
\crefname{originallemma}{original lemma}{original lemmas}
\Crefname{originallemma}{Original lemma}{Original lemmas}
\makeatletter
\newlabel{sites-cohomology-lemma-apply-Leray}{{0733}{source}{}{sourcepointer-5}{}}
\newlabel{sites-cohomology-lemma-apply-Leray@cref}{{[originallemma][0][]0733}{[1][0][]source}{}{}{}}
\makeatother
\title{Cohomology of a coherent module on a proper algebraic space over a Noetherian affine base is a finite module}
\author{The Stacks Project Authors}
\date{stacks-project:spaces-cohomology.tex:local-snapshot-20261012; GNU FDL 1.2 or later}
\begin{document}
\maketitle
\noindent\textit{Editorial note.} All delivered mathematical text is the author's own, copied from the pinned Stacks Project chapter file spaces-cohomology.tex; no statement or proof step was rewritten, reordered or removed.\par\medskip
\noindent\textit{Editorial note.} Preamble provenance: the pinned chapter file opens with \textbackslash{}input\{preamble\} and carries no preamble text of its own, so the project's single Stacks-wide preamble staged at sources/stacks-spaces-cohomology/preamble.tex is excerpted verbatim as three spans, lines 8-12 (verbatim package and the project's own reference/slogan/history display environments), lines 14-18 (Xy-pic setup) and lines 154-217 (theorem declarations and mathematical macros). The wrapper class and the amsmath/amssymb/amsthm/graphicx base packages are supplied by the assembler. No macro was invented or redefined.\par\medskip
\noindent\textit{Editorial note.} Same-chapter pointers are located at their \textbackslash{}label lines inside the pinned chapter file; cross-chapter pointers are located in the project's own tag list staged in the snapshot (sources/stacks-chow/tags.txt), at the verbatim lines carrying their canonical tags.\par\medskip
\noindent\textit{Allowed standard prerequisites (editorial).}\begin{itemize}
\item Coherence of the higher direct images of a coherent module under a proper morphism. Same chapter spaces-cohomology.tex, label lemma-proper-pushforward-coherent (invoked at proof line 4020). It is the affine case of this lemma that is being used: already knowing \$R\textasciicircum{}if\_*\textbackslash{}mathcal\{F\}\$ coherent gives the finiteness conclusion.
\item Quasi-coherence of the higher direct images along a quasi-compact quasi-separated morphism. Same chapter spaces-cohomology.tex, label lemma-higher-direct-image (invoked at proof line 4021). It makes \$R\textasciicircum{}if\_*\textbackslash{}mathcal\{F\}\$ a quasi-coherent module on \$\textbackslash{}Spec(A)\$, hence the module associated to its global sections.
\item The Leray identification of cohomology with global sections of a higher direct image, and vanishing on an affine scheme. Cohomology on Sites, label sites-cohomology-lemma-apply-Leray (invoked at proof line 4026) and Cohomology of Schemes, label coherent-lemma-quasi-coherent-affine-cohomology-zero (invoked at proof line 4028). Together they give \$H\textasciicircum{}i(X, \textbackslash{}mathcal\{F\}) = \textbackslash{}Gamma(\textbackslash{}Spec(A), R\textasciicircum{}if\_*\textbackslash{}mathcal\{F\})\$.
\item Coherence on a Noetherian affine base is finite generation of the global sections. Same chapter spaces-cohomology.tex, label lemma-coherent-Noetherian (invoked at proof line 4030). It turns the coherence of \$R\textasciicircum{}if\_*\textbackslash{}mathcal\{F\}\$ into finite generation of \$H\textasciicircum{}i(X, \textbackslash{}mathcal\{F\})\$ over \$A\$.
\end{itemize}\medskip
\ifdefined\hypertarget\hypertarget{sourcepointer-1}{}\fi
\noindent\textit{Source reference.} Original lemma 01XB: sources/stacks-chow/tags.txt, source lines 2512--2512. Navigation pointer for \textbackslash{}ref\{coherent-lemma-quasi-coherent-affine-cohomology-zero\} used inside the counted span (line 4029); the argument that consumes it is retained in the span.\par\medskip
\ifdefined\hypertarget\hypertarget{sourcepointer-2}{}\fi
\noindent\textit{Source reference.} Original lemma 07UB: sources/stacks-spaces-cohomology/pinned-original-source.tex, source lines 2094--2094. Navigation pointer for \textbackslash{}ref\{lemma-coherent-Noetherian\} used inside the counted span (line 4030); the argument that consumes it is retained in the span.\par\medskip
\ifdefined\hypertarget\hypertarget{sourcepointer-3}{}\fi
\noindent\textit{Source reference.} Original lemma 0720: sources/stacks-spaces-cohomology/pinned-original-source.tex, source lines 339--339. Navigation pointer for \textbackslash{}ref\{lemma-higher-direct-image\} used inside the counted span (line 4021); the argument that consumes it is retained in the span.\par\medskip
\ifdefined\hypertarget\hypertarget{sourcepointer-4}{}\fi
\noindent\textit{Source reference.} Original lemma 08AR: sources/stacks-spaces-cohomology/pinned-original-source.tex, source lines 3835--3835. Navigation pointer for \textbackslash{}ref\{lemma-proper-pushforward-coherent\} used inside the counted span (line 4020); the argument that consumes it is retained in the span.\par\medskip
\ifdefined\hypertarget\hypertarget{sourcepointer-5}{}\fi
\noindent\textit{Source reference.} Original lemma 0733: sources/stacks-chow/tags.txt, source lines 9047--9047. Navigation pointer for \textbackslash{}ref\{sites-cohomology-lemma-apply-Leray\} used inside the counted span (line 4026); the argument that consumes it is retained in the span.\par\medskip
\begin{lemma}
\label{lemma-proper-over-affine-cohomology-finite}
Let $A$ be a Noetherian ring.
Let $f : X \to \Spec(A)$ be a proper morphism of algebraic spaces.
Let $\mathcal{F}$ be a coherent $\mathcal{O}_X$-module.
Then $H^i(X, \mathcal{F})$ is finite $A$-module for all $i \geq 0$.
\end{lemma}

\begin{proof}
This is just the affine case of Lemma \ref{lemma-proper-pushforward-coherent}.
Namely, by Lemma \ref{lemma-higher-direct-image} we know that
$R^if_*\mathcal{F}$ is a quasi-coherent sheaf. Hence it is the quasi-coherent
sheaf associated to the $A$-module
$\Gamma(\Spec(A), R^if_*\mathcal{F}) = H^i(X, \mathcal{F})$.
The equality holds by
Cohomology on Sites, Lemma \ref{sites-cohomology-lemma-apply-Leray}
and vanishing of higher cohomology groups of quasi-coherent modules
on affine schemes (Cohomology of Schemes, Lemma
\ref{coherent-lemma-quasi-coherent-affine-cohomology-zero}).
By Lemma \ref{lemma-coherent-Noetherian} we see $R^if_*\mathcal{F}$ is
a coherent sheaf if and only if $H^i(X, \mathcal{F})$
is an $A$-module of finite type. Hence
Lemma \ref{lemma-proper-pushforward-coherent} gives us the conclusion.
\end{proof}

\end{document}
